\src{175}
\result{PROPOSITION}{4} \uline{Soit} $\symup{A}$ \uline{un anneau commutatif avec unité, soit}
$\symup{(x_1,\ldots,x_p)}$ \uline{une suite d'éléments de} $\symup{A}$ \uline{telle que pour} $\symup{1\leq i\leq p}$,
\uline{l'image de} $\symup{x_i}$ \uline{dans le quotient de} $\symup{A}$ \uline{par l'idéal engendré par}
$\symup{(x_1,\ldots,x_{i-1})}$ \uline{ne soit pas diviseur de} $\symup{0}$. \uline{Soit} $\symup{J}$ \uline{l'idéal engendré
par les} $\symup{x_i}$. \uline{Alors} $\symup{J/J^{2}}$ \uline{est un} $\symup{(A/J)}$\uline{-module libre ayant pour base les
images canoniques des} $\symup{x_i}$, \uline{le complexe} $\symup{K_{\underline{x}}}$ \uline{est une résolution
libre de} $\symup{A/J}$, \uline{et pour tout} $\symup{A}$\uline{-module} $\symup{M}$, \uline{les homomorphismes} (3.1) \uline{en
dimension} $\symup{p}$ \uline{sont bijectifs. Il en est de même des homomorphismes analogues
définis pour des degrés} $\symup{i}$ \uline{quelconques pourvu que} $\symup{J.M=0}$.

(Le point essentiel, dont tous les autres résultent, est l'acyclicité de
$\symup{K_{\underline{x}}}$, qui est un fait bien connu sous les conditions indiquées).

\result{COROLLAIRE}{1} $\symup{A}$ \uline{et} $\symup{J}$ \uline{étant comme ci-dessus, supposons de plus que}
$\symup{A}$ \uline{soit une algèbre affine régulière de} $\symup{\dim n}$ \uline{sur un corps parfait} $\symup{\underline{k}}$, \uline{et
que} $\symup{A/J}$ \uline{soit une algèbre affine régulière. Désignons par} $\symup{\Omega^{1}(A)}$,
$\symup{\Omega^{1}(A/J)}$ \uline{les modules de différentielles de Kähler. Alors on a un
isomorphisme canonique compatible avec la localisation},
\begin{equation}\tag{3.5}
  \symup{\operatorname{Ext}_{A}^{p}\bigl(\Omega^{n-p}(A/J),\Omega^{n}(A)\bigr)=A/J}
\end{equation}
En effet, $\symup{\Omega^{n-p}(A/J)}$ est un $\symup{(A/J)}$-module projectif de rang $\symup{1}$, de même
$\symup{\Omega^{n}(A)}$ est un $\symup{A}$-module projectif de rang $\symup{1}$,\footnote{Le texte imprimé dit « libre » pour les deux modules et « rang $\symup{n}$ » pour $\symup{\Omega^{n}(A)}$. Les modules de formes de degré maximal sont projectifs de rang $\symup{1}$ ; leur liberté globale ne découle pas des hypothèses. L'identification ci-dessous résulte de la suite conormale en prenant les déterminants.}
donc le premier membre est égal à
\[
  \symup{\operatorname{Ext}_{A}^{p}(A/J,A)
  \otimes\Omega^{n-p}(A/J)'\otimes\Omega^{n}(A)}
\]
(où le symbole \fquote{$\symup{\prime}$} désigne le $\symup{(A/J)}$-module dual). Le produit tensoriel des
deux derniers facteurs s'identifie à $\symup{\bigwedge^{p}(J/J^{2})}$, donc le tout
s'identifie à
$\symup{\operatorname{Ext}_{A}^{p}\!\left(A/J,\bigwedge^{p}(J/J^{2})\right)}$,
donc, en vertu de la proposition, à
\[
  \symup{\operatorname{Hom}_{A}\!\left(\bigwedge^{p}J/J^{2},
  \bigwedge^{p}J/J^{2}\right),}
\]
i.e. à $\symup{A/J}$. En particulier, il y a dans
$\symup{\operatorname{Ext}_{A}^{p}\bigl(\Omega^{n-p}(A/J),\Omega^{n}(A)\bigr)}$ un
élément privilégié, correspondant à l'élément unité de $\symup{A/J}$, appelé
\uline{classe fondamentale} de l'idéal $\symup{J}$ dans $\symup{A}$. (Elle peut en fait se
définir sous des conditions sensiblement plus larges). On peut mettre le
corollaire 1 sous une forme plus géométrique et plus globale :

\result{COROLLAIRE}{2} \uline{Soient} $\symup{X}$ \uline{une variété non singulière sur un corps
algébriquement clos} $\symup{\underline{k}}$, $\symup{Y}$ \uline{une sous-variété fermée non singulière de} $\symup{X}$,
$\symup{\Osh_X}$ \uline{le faisceau des anneaux locaux de} $\symup{X}$, $\symup{\Osh_Y}$ \uline{le faisceau des anneaux
locaux de} $\symup{Y}$ \uline{considéré comme un faisceau quotient de} $\symup{\Osh_X}$. \uline{Soient} $\symup{n}$ \uline{la
dimension de} $\symup{X}$, $\symup{n-p}$ \uline{celle de} $\symup{Y}$.

